%============================================================
% MTH 112 Project - Angles
% Updated 202504
%============================================================

\section{Angles} \label{section-angles}
In this section we will draw angles in standard position, convert between degrees and radians and degrees-minutes-seconds, find coterminal angles, and find the length of a circular arc.

\textbook{\href{https://openstax.org/books/algebra-and-trigonometry-2e/pages/7-1-angles}{7.1}}

\subsection*{Preparation Exercises} \label{preparation-exercises-section-angles}

Answer the following without using a calculator. These questions are intended to check the prerequisite skills needed to complete the rest of the lab.

\begin{multicols}{2}
	\begin{myPrep}
		Solve the equation $T=\frac{v}{w^2}$ for $w$.
	\end{myPrep}
	\vfill
	
	\begin{myPrep}
		Simplify the expression $80\cdot\frac{\pi}{180}$ without a calculator. Reduce to simplest form.
	\end{myPrep}
	\vfill
\end{multicols}
\vfill

\begin{myPrep}
In terms of time, how many minutes total are in 3 hours and 53 minutes?
\end{myPrep}
\vfill

\begin{myPrep}
	\begin{minipage}[t]{0.5\linewidth}
		Draw a line segment from the point $(5,-3)$ to the point $(-1,2)$. Find the distance between those two points using the Pythagorean Theorem. Leave your answer in exact form.
		\vfill
	\end{minipage}
	\begin{minipage}[t]{0.5\linewidth}
		\begin{center}\captionsetup{type=figure} \captionof{figure}{A grid to plot your points}
			\begin{tikzpicture}
				\begin{axis}[
					framed,
					xmin=-6,xmax=6,
					ymin=-6,ymax=6,
			       	xtick={-4,-2,2,4},
			       	minor xtick={-5,-3,-1,1,3,5},
			       	ytick={-4,-2,2,4},
			       	minor ytick={-5,-3,-1,1,3,5},
			       	grid=both
					]
				\end{axis}
			\end{tikzpicture}\label{fig:section-angles-myprep1}
		\end{center}
	\end{minipage}
\end{myPrep}
\vfill

%======================================================
\newpage
%======================================================

\subsection*{Practice Exercises} \label{practice-exercises-section-angles}

\begin{myPractice}
	Find an angle that is coterminal to $\ang{-40}$ and is between $\ang{360}$ and $\ang{720}$.
		% Solution: Since $\ang{-40}$ is in Quadrant IV, $\ang{40}$ below the $x$-axis, a coterminal angle would be $\ang{360}-\ang{40}=\ang{320}$. However, this isn't between $\ang{360}$ and $\ang{720}$, so we can add another $\ang{360}$ and we get $\ang{320}+\ang{360}=\ang{680}$. Another way to find this angle is just to subtract $\ang{40}$ from $\ang{720}$: $\ang{720}-\ang{40}=\ang{680}$.
		\vfill
\end{myPractice}

\begin{myPractice}
	Match the pairs of coterminal angles.
			\begin{multicols}{2}
				\begin{itemize}
					\item \ang{30}
					\item \ang{270}
					\item \ang{-220}
					\item \ang{-60}
					\item \ang{0}
					\item \ang{150}
					\columnbreak
					\item \ang{720}
					\item \ang{-90}
					\item \ang{390}
					\item \ang{510}
					\item \ang{-420}
					\item \ang{500}
				\end{itemize}
			\end{multicols}
\end{myPractice}

\begin{myPractice}
	Convert the angles.
			\begin{enumerate}
				\item Convert $\ang{28}$ to radians. Leave your answer as a fraction in terms of $\pi$.
				\vfill
				\item Convert $\frac{7\pi}{9}$ radians to degrees. Round to 6 digits behind the decimal place.
				\vfill

%======================================================
\newpage
%======================================================

				\item Convert $\ang{14;20;36}$ to a decimal measurement of degrees. 
				\vfill
				\item Convert $\ang{76.85}$ to degrees-minutes-seconds.
				\vfill
			\end{enumerate}
\end{myPractice}

\begin{myPractice}
	Use the formula $s=r\theta$ to find the missing values.
			\begin{enumerate}
				\item Find the arc length of on a circle of radius $6$ with an interior angle of $\frac{4\pi}{7}$.
				\vfill
				\item Find the interior angle of an arc length on a circle of radius $12$ with an arc length of $\frac{3\pi}{2}$.
				\vfill
				\item Find the arc length of on a circle of radius $10\pi$ with an interior angle of $\ang{20}$.
				\vfill
			\end{enumerate}
		\vfill
\end{myPractice}

%======================================================
\newpage
%======================================================

\subsection*{Definitions} \label{definitions-section-angles}
\begin{multicols}{2}
	\begin{myDefinition}[{Angle}]~\\[0.5mm]
		An \textbf{angle} is formed when two lines intersect at a point. We can measure the \say{size} of this angle in degrees or radians.
	\end{myDefinition}
	\begin{myDefinition}[{Ray}]~\\[0.5mm]
		A \textbf{ray} is a portion of a line that starts at a point and extends infinitely in one direction.
	\end{myDefinition}
	\begin{myDefinition}[{Degree}]~\\[0.5mm]
		A \textbf{degree} is an angle measurement equal to $\frac{1}{360}$ of a full rotation.
	\end{myDefinition}
	\begin{myDefinition}[{Right Angle}]~\\[0.5mm]
		 A \textbf{right angle} is an angle of $\ang{90}$.
	\end{myDefinition}
	\begin{myDefinition}[{Radian}]~\\[0.5mm]
			A \textbf{radian} is a unit of measurement of angle size, and 1 radian is defined to be the angle that is made when the arc length for that angle on a circle of radius 1 is equal to 1. \href{http://tiny.cc/radian-definition}{Check out this Desmos link} and click the "play" button on the slider for a visual representation of this concept. \qrcode[height=1cm]{http://tiny.cc/radian-definition}
	\end{myDefinition}

	\begin{center}
		\captionsetup{type=figure} \captionof{figure}{Standard Angles}
		\begin{tikzpicture}
			\begin{axis}[
					xmin=-1,xmax=1,
					ymin=-1,ymax=1,
		       		xtick={-1,1},
		       		ytick={-1,1},
		       		clip=false,
					]
					\draw 	(axis cs:0,0) circle [radius=1];
					\addplot[mark=*,color=first] coordinates{(1,0)} node[below left] {\footnotesize$\theta=0$};

					\draw   [color=second,opacity=0.2](0,0)--(0.866,0.5);
					\addplot[mark=*,color=second] coordinates{(0.866,0.5)} node[left,rotate=30] {\footnotesize$\theta=\frac{\pi}{6}$};
					\draw   [color=third,opacity=0.2](0,0)--(0.707,0.707);
					\addplot[mark=*,color=third] coordinates{(0.707,0.707)} node[left,rotate=45] {\footnotesize$\theta=\frac{\pi}{4}$};

					\draw   [color=fourth,opacity=0.2](0,0)--(0.5,0.866);
					\addplot[mark=*,color=fourth] coordinates{(0.5,0.866)} node[left,rotate=60] {\footnotesize$\theta=\frac{\pi}{3}$};

					\addplot[mark=*,color=fifth] coordinates{(0,1)}node[below right,rotate=-90] {\footnotesize$\theta=\frac{\pi}{2}$};

					\draw   [color=fourth,opacity=0.2](0,0)--(-0.5,0.866);
					\addplot[mark=*,color=fourth] coordinates{(-0.5,0.866)}node[right,rotate=-60] {\footnotesize$\theta=\frac{2\pi}{3}$};

					\draw   [color=third,opacity=0.2](0,0)--(-0.707,0.707);
					\addplot[mark=*,color=third] coordinates{(-0.707,0.707)}node[right,rotate=-45] {\footnotesize$\theta=\frac{3\pi}{4}$};

					\draw   [color=second,opacity=0.2](0,0)--(-0.866,0.5);
					\addplot[mark=*,color=second] coordinates{(-0.866,0.5)}node[right,rotate=-30] {\footnotesize$\theta=\frac{5\pi}{6}$};

					\addplot[mark=*,color=first] coordinates{(-1,0)} node[above right] {\footnotesize$\theta=\pi$};
					\draw   [color=second,opacity=0.2](0,0)--(-0.866,-0.5);
					\addplot[mark=*,color=second] coordinates{(-0.866,-0.5)} node[right,rotate=30] {\footnotesize$\theta=\frac{7\pi}{6}$};
					\draw   [color=third,opacity=0.2](0,0)--(-0.707,-0.707);
					\addplot[mark=*,color=third] coordinates{(-0.707,-0.707)} node[right,rotate=45] {\footnotesize$\theta=\frac{5\pi}{4}$};
					\draw   [color=fourth,opacity=0.2](0,0)--(-0.5,-0.866);
					\addplot[mark=*,color=fourth] coordinates{(-0.5,-0.866)} node[right,rotate=60] {\footnotesize$\theta=\frac{4\pi}{3}$};

					\addplot[mark=*,color=fifth] coordinates{(0,-1)} node[above left,rotate=-90] {\footnotesize$\theta=\frac{3\pi}{2}$};
					\draw   [color=second,opacity=0.2](0,0)--(0.866,-0.5);
					\addplot[mark=*,color=second] coordinates{(0.866,-0.5)} node[left,rotate=-30] {\footnotesize$\theta=\frac{11\pi}{6}$};
					\draw   [color=third,opacity=0.2](0,0)--(0.707,-0.707);
					\addplot[mark=*,color=third] coordinates{(0.707,-0.707)} node[left,rotate=-45] {\footnotesize$\theta=\frac{7\pi}{4}$};
					\draw   [color=fourth,opacity=0.2](0,0)--(0.5,-0.866);
					\addplot[mark=*,color=fourth] coordinates{(0.5,-0.866)} node[left,rotate=-60] {\footnotesize$\theta=\frac{5\pi}{3}$};
			\end{axis}
		\end{tikzpicture}
		\label{fig:Unit-circle-values-angles}
	\end{center}
% \end{multicols}

% \begin{multicols}{2}
	\begin{myDefinition}[{Initial and Terminal sides}]~\\[0.5mm]
		An angle is created by imagining two rays originating from the same point: the \textbf{initial side} of the angle is stationary and the \textbf{terminal side} is rotated around until the desired angle is created. Angles measured from initial to terminal side in a counter-clockwise direction are positive. Angles measured from initial to terminal side clockwise are negative.
	\end{myDefinition}
	\begin{myDefinition}[{Standard Position}]~\\[0.5mm]
		An angle is said to be in \textbf{standard position} if the initial and terminal sides of the angle meet at the origin and the initial side of the angle is along the positive $x$-axis.
	\end{myDefinition}
	\begin{myDefinition}[{Coterminal}]~\\[0.5mm]
		Two angles are \textbf{coterminal} if their initial and terminal sides align, but they differ by an integer number of full rotations. In standard position, this means that the two angles are \say{in the same place} are different by some integer multiple of $\ang{360}$ (or $2\pi$ radians).
	\end{myDefinition}
	\begin{myDefinition}[{Arc length}]~\\[0.5mm]
		An \textbf{arc length} is a distance along a part of the circumference of a circle. A formula for arc length is $s=r\theta$, where $s$ stands for arc length along a circle of radius $r$, with an interior angle $\theta$ centered at the center of the circle.
	\end{myDefinition}
	
	\begin{center}
		\captionsetup{type=figure} \captionof{figure}{Coterminal angles $\theta$ and $\alpha$}
		\begin{tikzpicture}
			\begin{axis}[
				% framed,
				xmin=-10,xmax=10,
				ymin=-10,ymax=10,
		       	xtick=\empty,
		       	% minor xtick={-5,5},
		       	ytick=\empty,
		       	% minor ytick={-5,5},
		       	% grid=both,
		       	clip=false
				]
				\draw 	(axis cs:0,0) circle [radius=10];
				
				\draw   [color=first, very thick](10,0)--(5,0)coordinate[label=below:\tiny{initial side}]--(0,0)--(-6,8) coordinate[label=above left:{$\theta$ and $\alpha$ are coterminal}];

				\draw 	(-3,4) node[below,color=third,rotate=-atan(4/3)] {\tiny{terminal side}};
				
				\addplot[variable=\t, domain=0:(atan(-4/3)+180), samples=50, color=first, very thick,-{latex}]({2*cos(t)}, {2*sin(t)});
				
				\draw 	(2,1) node[right,color=first] {$\theta$};
				
				% \addplot[variable=\t, domain=0:(atan(-4/3)+180), samples=50, color=second,very thick]({10*cos(t)}, {10*sin(t)});
				
				\draw[postaction={decorate},very thick,color=second,decoration={text along path,raise=5pt,text color=second,text={arc length for angle {$\theta$}},text align=center}] (-6,8) arc [start angle=atan(-4/3)+180,end angle=0,radius=10];
				
				\addplot[variable=\t, domain=(atan(-4/3)+180):360, samples=50, color=fourth,{latex}-]({2*cos(t)}, {2*sin(t)});
				
				\draw 	(-2,-1) node[left,color=fourth] {$\alpha$};
				\end{axis}
		\end{tikzpicture}
		\label{fig:angle-defs}
	\end{center}
\end{multicols}

%======================================================
\newpage
%======================================================

\subsection*{Exit Exercises} \label{exit-exercises-section-angles}

\begin{myExit} 
Explain why $\ang{-90}$ is not {\it the same} as $\ang{270}$.
\end{myExit}
\vfill

\begin{myExit} 
Convert $\ang{34;29;24}$ to decimal degrees.
\end{myExit}
\vfill
\vfill

\begin{myExit} 
	Convert the angle $\ang{144}$ to radians. Write you answer in exact form as a fraction with $\pi$ in it.
\end{myExit}
\vfill

\begin{myExit} 
	\begin{minipage}[t]{0.6\linewidth}
		How big of an arc length will an angle of $\frac{3\pi}{8}$ radians create on a circle of radius 6? Draw this angle and arc length on the circle of radius 6 provided.
	\end{minipage}
	\begin{minipage}[t]{0.4\linewidth}
		\begin{center}
			\captionsetup{type=figure} \captionof{figure}{A circle of radius 6}
			\begin{tikzpicture}
				\begin{axis}[
					framed,
					xmin=-8,xmax=8,
					ymin=-8,ymax=8,
	        		xtick={-6,-4,-2,2,4,6},
	        		ytick={-6,-4,-2,2,4,6},
	        		grid=both
					]
					\draw[thick,color=first] (0,0) circle (6);
				\end{axis}
			\end{tikzpicture}
		\end{center}
		\label{fig:section-angles-myexit-circle}
	\end{minipage}
\end{myExit}
\vfill
\exitlikert{angles}